\documentclass[a4paper,12pt]{article}

% Language and encoding
\usepackage[ngerman]{babel}

% Page layout
\usepackage{geometry}
\geometry{a4paper, top=30mm, left=25mm, right=25mm, bottom=30mm}

% Fonts and spacing
\usepackage{setspace}
\onehalfspacing
\usepackage{lmodern}

% Header and footer
\usepackage{fancyhdr}
\pagestyle{fancy}
\fancyhf{}
\fancyhead[L]{\leftmark}
\fancyhead[R]{\thepage}

% Math and symbols
\usepackage{amsmath, amssymb}

% Graphics
\usepackage{graphicx}
\usepackage{float}
\usepackage{caption}
\usepackage{subcaption}

% Tables
\usepackage{booktabs}
\usepackage{siunitx}
\sisetup{locale = DE} % Use German number formatting

% Links
\usepackage[hidelinks]{hyperref}

% Title page
\title{
    \Huge{Einführung in die Simulation\\Maschinenfehler} \\
    \vspace{0.5cm}
    \large{Paris Lodron Universität Salzburg}
}
\author{
    Schmidt Linda, Weilert Alexander, Schlager Andreas \\
}
\date{\today}

% Main document
\begin{document}

    \maketitle
    \thispagestyle{empty}
    \newpage

    \tableofcontents
    \newpage

    \section{Einleitung}
    Hier wird eine Einführung in das Thema des Berichts gegeben. Zum Beispiel, welches System simuliert wird und warum.

    \newpage

    \section{Modellierung}
    Hier wird das gewählte Simulationsmodell vorgestellt. Dazu gehört:
    \begin{itemize}
        \item Annahmen
        \item Modellstruktur (eventuell mit Diagrammen)
        \item Formale Beschreibung (z.\,B. Zustände, Ereignisse, Wahrscheinlichkeiten)
    \end{itemize}

    \newpage

    \section{Implementierung}
    Details zur technischen Umsetzung (Programmiersprache, Simulationsumgebung, Bibliotheken).

    \newpage

    \section{Ergebnisse}
    Darstellung und Auswertung der durchgeführten Simulationsexperimente:
    \begin{itemize}
        \item Simulationsparameter
        \item Ergebnisse in Tabellen und Diagrammen
        \item Statistische Auswertungen
    \end{itemize}

    \section{Fazit}
    Zusammenfassung der wichtigsten Erkenntnisse aus dem Bericht.

    \newpage
    \appendix

    \section{Anhang}
    Zusätzliche Informationen, z.\,B. Quellcode-Auszüge oder vollständige Ergebnislisten.
\end{document}
